% \makebibliography

% \renewcommand{\refname}{\centering{\large Мой собственный заголовок}}

\addcontentsline{toc}{chapter}{\textbf{Список использованных источников}}

\begin{thebibliography}{99}

\bibitem{aho2008compilers}
Ахо А., Сети Р., Ульман Д.
Компиляторы: принципы, технологии и инструментарий. – М.: ООО «И.Д. Вильямс», 2008. – 1184 с.

\bibitem{serebryakov_construct}
Серебряков В. А., Галочкин М. П.
Основы конструирования компиляторов.

\bibitem{aho1978syntax}
Ахо А., Ульман Дж.
Теория синтаксического анализа, перевода и компиляции. Том 1. – М.: «Мир», 1978.

\bibitem{parr2013antlr}
Parr T. The Definitive ANTLR4 Reference. – M.: Pragmatic Bookshelf, 2013.

\bibitem{lattner2012llvm}
Chris Lattner.  
The Design of LLVM // Dr. Dobb’s Journal. – 2012.

\end{thebibliography}